%% Diagramme des cas d'utilisation (uc) du pont transbordeur de Nantes.
%% Frontière du système = rectangle bleu ; acteurs à l'extérieur ; relations
%% «extend» (fonction non indispensable) et «include» (fonction indispensable).
\documentclass[tikz,border=4pt]{standalone}
\usepackage{liaisons}
\begin{document}
\begin{tikzpicture}[x=1cm,y=1cm,
  cadre/.style={draw=black!70, line width=0.9pt},
  trait/.style={line width=1pt, line cap=round, line join=round},
  cas/.style={draw=solideD, line width=1pt, fill=solideD!12},
  txt/.style={font=\scriptsize, align=center, inner sep=1pt,
              execute at begin node={\hyphenpenalty=10000\exhyphenpenalty=10000}},
  lien/.style={line width=0.8pt, black!75},
  pointille/.style={-{Straight Barb[length=4pt,width=6pt]}, line width=0.9pt,
                    dash pattern=on 3pt off 2pt, solideD},
  ster/.style={font=\scriptsize, inner sep=2pt, anchor=west}]

%% ---- macros de pictogrammes -----------------------------------------
\newcommand\acteur[3]{%
  \begin{scope}[shift={(#1,#2)}, trait, draw=#3]
    \draw (0,0.42) circle (0.15); \draw (0,0.27) -- (0,-0.12);
    \draw (-0.28,0.12) -- (0.28,0.12);
    \draw (-0.22,-0.5) -- (0,-0.12) -- (0.22,-0.5);
  \end{scope}}
\newcommand\cube[3]{%
  \begin{scope}[shift={(#1,#2)}, trait, draw=#3]
    \draw[fill=#3!12] (-0.3,-0.3) rectangle (0.2,0.2);
    \draw (-0.3,0.2) -- (-0.14,0.36) -- (0.36,0.36) -- (0.2,0.2);
    \draw (0.36,0.36) -- (0.36,-0.14) -- (0.2,-0.3);
  \end{scope}}

%% ================== cadre à onglet ====================================
\draw[cadre] (0,0) rectangle (9.4,5.5);
\draw[cadre] (0,5.05) -- (3.55,5.05) -- (3.8,5.28) -- (3.8,5.5);
\node[font=\scriptsize, anchor=west, inner sep=3pt] at (0.1,5.27)
     {\texttt{uc} [Paquet] Utilisation principale};

%% ================== frontière du système ==============================
\draw[draw=solideB, line width=1pt, fill=solideB!7] (2.55,0.25) rectangle (6.85,4.75);
\node[txt, text=solideB] at (4.7,4.5) {«useCaseModel»};
\node[font=\small\bfseries, text=solideB] at (4.7,4.15) {Pont transbordeur};

%% ================== les trois cas d'utilisation =======================
\draw[cas] (4.7,3.35) ellipse (1.85 and 0.48);
\node[txt, text width=3.3cm] at (4.7,3.35) {Déambuler sur la rue aérienne};
\draw[cas] (4.7,2.05) ellipse (1.85 and 0.48);
\node[txt, text width=3.3cm] at (4.7,2.05) {Franchir la Loire};
\draw[cas] (4.7,0.85) ellipse (1.85 and 0.48);
\node[txt, text width=3.3cm] at (4.7,0.85) {Permettre le passage de bateaux};

%% ================== relations «extend» et «include» ===================
\draw[pointille] (4.2,2.87) -- (4.2,2.53);
\node[ster, text=solideD] at (4.3,2.7) {«extend»};
\draw[pointille] (4.2,1.57) -- (4.2,1.33);
\node[ster, text=solideD] at (4.3,1.45) {«include»};

%% ================== acteurs principaux (à gauche) =====================
\acteur{1.15}{3.5}{solideB}
\node[txt, text width=1.9cm] at (1.15,2.75) {Utilisateur};
\draw[lien] (1.4,3.25) -- (3.05,2.35);
\cube{1.2}{1.35}{solideE}
\node[txt, text width=1.9cm] at (1.2,0.75) {Busway};
\draw[lien] (1.6,1.45) -- (3.05,1.8);

%% ================== acteurs secondaires (à droite) ====================
\cube{8.4}{3.6}{solideE}
\node[txt, text width=1.9cm] at (8.45,3.0) {La Loire};
\draw[lien] (8.15,3.4) -- (6.4,2.35);
\cube{8.4}{1.35}{solideE}
\node[txt, text width=1.9cm] at (8.45,0.75) {Bateaux};
\draw[lien] (8.15,1.25) -- (6.4,1.0);
\end{tikzpicture}
\end{document}
